\section{Introduction}
\label{sec:intro}

% --- motivation ---
Modern LLM agents persist context across sessions through external memory
systems such as Mem0~\citep{mem0}, A-MEM~\citep{amem}, and
MemGPT~\citep{memgpt2023}.  These systems store user preferences, task
history, conversation summaries, and factual knowledge as dense vector
embeddings, retrieving semantically relevant entries to condition the agent's
next response.  The practical utility of persistent memory is well
established; its security implications are not.

% --- the attack surface ---
An adversary who can influence the contents of an agent's memory---whether
through direct write access, crafted user interaction, or exploitation of
the agent's self-storage mechanism---can covertly redirect agent behavior
across arbitrary future sessions.  Unlike jailbreaks, which operate on a
single interaction, memory poisoning attacks \emph{persist}: a single
adversarial entry may remain in memory for weeks and trigger on every
related query.  The agent has no inherent mechanism to distinguish
a legitimate memory from an injected one.

% --- three attacks ---
Three recent papers have characterized different threat models for memory
poisoning.  \agentpoison{} \citep{chen2024agentpoison} requires write access
to the memory/knowledge base and optimizes a short trigger token sequence
such that any user query containing the trigger retrieves the adversarial
entry with high probability.  \minja{} \citep{dong2025minja} requires no
memory access: the attacker poses as a regular user and submits crafted
queries that cause the agent to auto-generate and store adversarial reasoning
chains.  \injecmem{} \citep{injecmem2026} achieves injection in a single
interaction via a ``retriever-agnostic anchor'' passage designed to
match many query types simultaneously.

% --- our contributions ---
This paper makes four contributions:

\begin{enumerate}
  \item \textbf{Unified evaluation framework.}  We implement all three attacks
  over a realistic FAISS-backed vector memory system with a 200-entry synthetic
  agent corpus and three paper-faithful success metrics: $\asrr$ (retrieval
  success rate), $\asra$ (action execution rate given retrieval), and $\asrt$
  (end-to-end task hijacking rate) as defined in \citet{chen2024agentpoison}.

  \item \textbf{Evaluation protocol correction.}  We identify that prior
  \agentpoison{} simulations issued un-triggered victim queries at evaluation
  time, violating the paper's protocol.  With the correct triggered-query
  evaluation and a centroid-targeting adversarial passage, $\asrr$ increases
  from 0.25 to 1.00 (Section~\ref{sec:experiments}).

  \item \textbf{Semantic Anomaly Detection (\sad{}).}  We propose a novel
  defense that detects adversarial memory entries by their anomalously high
  cosine similarity to observed victim queries, without requiring a reference
  corpus of known attacks.  \sad{} achieves $\tpr > 0.70$ with $\fpr < 0.10$
  across all three attack types.

  \item \textbf{Full attack-defense interaction matrix.}  We evaluate all
  $3 \times 5$ (attack, defense) pairs under a pre-ingestion filtering model,
  revealing that no single defense dominates uniformly and motivating composite
  defense portfolios.
\end{enumerate}

% --- paper organization ---
Section~\ref{sec:related} reviews related work.
Section~\ref{sec:threat} formalizes the threat model.
Sections~\ref{sec:attacks} and~\ref{sec:defenses} describe attack and defense
implementations.  Section~\ref{sec:experiments} reports empirical results.
Section~\ref{sec:discussion} discusses limitations and future work.
